\subsection{同调维数为1的环}

命题 7. --- 下列条件是等价的：

\begin{enumerate}
\def\labelenumi{(\roman{enumi})}
\item
  \(\mathrm{dh}(\mathbf{A})\leqslant 1\)
\item
  投射模的每个子模都是投射的，
\end{enumerate}

(ii') A 的每个理想都是投射的，

\begin{enumerate}
\def\labelenumi{(\roman{enumi})}
\setcounter{enumi}{2}
\item
  内射 A-模的每个商都是内射的，
\item
  对每个投射复形 C，存在一个拟同构 \(\varphi : \mathbf{C} \to \mathbf{H}(\mathbf{C})\) 使得 \(\mathbf{H}(\varphi) = 1_{\mathbf{H}(\mathbf{C})}\) ,
\item
  对每个内射复形 C，存在一个拟同构 \(\psi : \mathrm{H}(\mathbf{C}) \to \mathbf{C}\) 使得 \(\mathrm{H}(\psi) = 1_{\mathrm{H}(\mathbf{C})}\) .
\item
  (i) \(\Leftrightarrow\) (ii) \(\Leftrightarrow\) (iii)：这由 X，第 138 页，命题 4 得出。
\item
  \(\Rightarrow\) (iv)：设 C 为投射复形。若(ii)成立，则 C 的子模 B(C) 是投射的，从而由 X，第 35 页，注 b) 得(iv)。
\item
  \(\Rightarrow\) (v)：设 C 为内射复形。若(iii)成立，则商 \(B_{n}(C)\) （\(C_{n+1}\) 的）对每个 \(n\) 都是内射的，从而由 X，第 35 页，注 a) 得(v)。
\item
  \(\Rightarrow\) (ii)：设 P 为投射 A-模，M 为 P 的子模， \(i: \mathbf{M} \to \mathbf{P}\) 为典范单射。设 \(p: \mathbf{L} \to \mathbf{M}\) 为自由模 L 到 M 上的满同态。考虑投射复形 C，使得 \(\mathbf{C}_1 = \mathbf{L}\) , \(\mathbf{C}_0 = \mathbf{P}\) , \(\mathbf{C}_i = 0\) （当 \(i \neq 0\) , \(1\)）， \(d_1 = i \circ p\) . 若(iv)满足，设 \(\varphi: \mathbf{C} \to \mathbf{H}(\mathbf{C})\) 为拟同构使得 \(\mathbf{H}(\varphi) = 1_{\mathbf{H}(\mathbf{C})}\) 。由于 \(\mathbf{H}_1(\mathbf{C}) = \operatorname{Ker} p\) , \(\varphi_1\) 是 L 到 \(\operatorname{Ker} p\) 上的投影算子，正合序列
\end{enumerate}

\[
0 \rightarrow \operatorname{Ker} p \rightarrow \mathrm{L} \xrightarrow {p} \mathrm{M} \rightarrow 0
\]

是分裂的，且 M 同构于 L 的一个直和项，因此是投射的。

\begin{enumerate}
\def\labelenumi{(\alph{enumi})}
\setcounter{enumi}{21}
\tightlist
\item
  \(\Rightarrow\) (iii)：设 I 为内射模，M 为 I 的商， \(\pi : \mathrm{I} \to \mathrm{M}\) 为典范投影。设 \(i: \mathrm{M} \to \mathrm{J}\) 为 M 到内射模 J 内的单同态。考虑内射复形 C，使得 \(C^0 = I\) , \(C^1 = J\) , \(C^i = 0\) （当 \(i \neq 0, 1\)）， \(d^0 = i \circ \pi\) 。若(v)满足，设 \(\psi : \mathrm{H}(\mathbf{C}) \to \mathbf{C}\) 为拟同构使得 \(\mathrm{H}(\psi) = 1_{\mathrm{H}(\mathbf{C})}\) 。由于
\end{enumerate}

\[
\mathrm{H} ^ {1} (\mathrm{C}) = \text { Coker } i  ,
\]

\(\psi^{1}\) 是典范投影 J → Coker i 的一个截面，且 M 是 J 的一个直和项，因此是内射的。

\begin{enumerate}
\def\labelenumi{(\roman{enumi})}
\setcounter{enumi}{1}
\tightlist
\item
  \(\Rightarrow\) (ii')：这是平凡的。
\end{enumerate}

(ii') \(\Rightarrow\) (ii)：这由 VII，第 3 节，定理 1 的推论 1 得出。

例. --- 若 A 是主理想环，则 dh(A) ≤ 1。

注. --- 若 A 是（交换）整环，则前述条件也等价于：

满足这些条件的整环称为戴德金环（参见 X，第 204 页，习题 12 以及 AC，VII，第 2 节，第 2 号，定理 1）。

推论. --- 设 A 为同调维数 ≤ 1 的环，C 为投射 A-模的复形， \(\tilde{C}\) 为投射右 A-模的复形，C' 为内射 A-模的复形，P 为右 A-模，M 为左 A-模，n 为整数。则我们有分裂正合序列

\[
\begin{array}{r l} & 0 \to \bigoplus_ {p + q = n} \mathrm{H} _ {p} (\tilde {\mathrm{C}}) \otimes_ {\mathrm{A}} \mathrm{H} _ {q} (\mathrm{C}) \xrightarrow {\gamma} \mathrm{H} _ {n} (\tilde {\mathrm{C}} \otimes_ {\mathrm{A}} \mathrm{C}) \xrightarrow {\alpha} \bigoplus_ {p + q = n - 1} \mathrm{Tor} _ {1} ^ {\mathrm{A}} \left(\mathrm{H} _ {p} (\tilde {\mathrm{C}}), \mathrm{H} _ {q} (\mathrm{C})\right) \to 0, \\ & 0 \to \prod_ {p} \mathrm{Ext} _ {\mathrm{A}} ^ {1} \left(\mathrm{H} _ {p} (\mathrm{C}), \mathrm{H} ^ {n - p - 1} (\mathrm{C} ^ {\prime})\right) \xrightarrow {\beta} \mathrm{H} ^ {n} (\mathrm{Homgr} _ {\mathrm{A}} (\mathrm{C}, \mathrm{C} ^ {\prime})) \\ & \qquad \qquad \qquad \xrightarrow {\lambda_ {p}} \prod_ {p} \mathrm{Homgr} _ {\mathrm{A}} \left(\mathrm{H} _ {p} (\mathrm{C}), \mathrm{H} ^ {n - p} (\mathrm{C} ^ {\prime})\right) \to 0, \\ & 0 \to \mathrm{P} \otimes_ {\mathrm{A}} \mathrm{H} _ {n} (\mathrm{C}) \xrightarrow {\gamma_ {p}} \mathrm{H} _ {n} (\mathrm{P} \otimes_ {\mathrm{A}} \mathrm{C}) \xrightarrow {\alpha_ {p}} \mathrm{Tor} _ {1} ^ {\mathrm{A}} \left(\mathrm{P}, \mathrm{H} _ {n - 1} (\mathrm{C})\right) \to 0, \\ & 0 \to \mathrm{Ext} _ {\mathrm{A}} ^ {1} \left(\mathrm{H} _ {n - 1} (\mathrm{C}), \mathrm{M}\right) \xrightarrow {\beta_ {p}} \mathrm{H} ^ {n} (\mathrm{Homgr} _ {\mathrm{A}} (\mathrm{C}, \mathrm{M})) \xrightarrow {\lambda_ {p}} \mathrm{Hom} _ {\mathrm{A}} \left(\mathrm{H} _ {n} (\mathrm{C}), \mathrm{M}\right) \to 0. \end{array}
\]

由于 \(Z(C)\) , \(B(C)\) , \(Z(\tilde{C})\) 与 \(B(\tilde{C})\) 是投射的且 \(B(C')\) 是内射的，这由 \(X\) ，第 78 页，推论 2，第 96 页，推论 1 和第 98 页，推论 2 得出。

